\section{Introduction}\label{sec:intro}

\subsection{Results}\label{ssec:results}
Let $\FQ$ be the set of primitive Dirichlet characters $\chi$ of modulus $q$ with $1<q\le Q$. For $\chi\in\FQ$ and $0\le T_1<T_2$ let
$N_\chi(T_1,T_2)$ be the number of zeros $\rho=\beta+i\gamma$ of $L(s,\chi)$ with $0<\beta<1$ and $T_1<\gamma\le T_2$, counted with
multiplicity, and let $N^{s}_{0,\chi}(T_1,T_2)$ be the number of those that are simple and have $\beta=\frac12$. Throughout, $r_0\ge3$ and
$\eps>0$ are fixed and
\[
T=T(Q)=(\log Q)^{r_0+\eps}.
\]
(The height exponent is called $r$ in~\cite{DSouza} and in the Lean statements; here $r$ denotes a modulus, see the notation table
in \S\ref{ssec:notation}.) The preprint~\cite{DSouza} proved that $\sum_{\chi\in\FQ}N^s_{0,\chi}(T,2T)\ge(0.7212-o(1))\sum_{\chi\in\FQ}N_\chi(T,2T)$.
This paper proves the following two theorems.

\begin{theorem}[killed kernel]\label{thm:one-prime}
There are $Q_0$ and $c>0$, depending on $r_0$ and $\eps$, such that for every integer $Q\ge Q_0$
\[
\sum_{\chi\in\FQ}N^s_{0,\chi}(T,2T)\ \ge\ \Bigl(0.7235-c\,\frac{\log\log Q}{\log Q}\Bigr)\sum_{\chi\in\FQ}N_\chi(T,2T).
\]
\end{theorem}

\noindent\emph{Status.} Proved on paper (\S\ref{sec:lemmaK}), with the corrections to the working draft listed in Appendix~\ref{app:corr}. In
Lean~4~\cite{Lean4} the statement is \texttt{ZetaShell.LemmaK.theorem\_one\_prime\_design}, a theorem with the three standard axioms and no
hypothesis beyond $r_0\ge3$, $\eps>0$, in the library \texttt{ZetaShell} of the accompanying Lean repository (\S\ref{ssec:formal}) \tagL.

\begin{theorem}[the Shell kernel]\label{thm:shell}
Assume the zero-density estimates \textup{(J)} and \textup{(M)} of \S\ref{ssec:shell-density}, which are theorems of
Jutila~\cite[Theorem~1]{Jutila77} and of Montgomery~\cite{Montgomery69} (in the form proved in~\cite[Th\'eor\`eme~20]{BombieriAst18}).
Let $0<\theta<503/1994$. There are $Q_0$ and $c>0$, depending on $r_0$, $\eps$ and $\theta$, such that for every integer $Q\ge Q_0$
\[
\sum_{\chi\in\FQ}N^s_{0,\chi}(T,2T)\ \ge\ \bigl(0.9059137927-c\,(\log Q)^{-\theta}\bigr)\sum_{\chi\in\FQ}N_\chi(T,2T).
\]
\end{theorem}

\noindent\emph{Status.} Proved on paper (\S\S\ref{sec:shell}--\ref{sec:proof}), with the corrections listed in Appendix~\ref{app:corr}. In Lean
(\S\ref{ssec:formal}) its statement of record, \texttt{ZetaShell.shell\_S53\_qle\_reduced}, carries (J) and (M) as its only displayed hypothesis. It is
a theorem in the library \texttt{ZetaShell} of the release: it depends on the three standard axioms only, and the census of the release finds no open
statement below it \tagL\ (Appendix~\ref{app:lean}). What is claimed can be read in two trusted files that depend on Mathlib~\cite{Mathlib} alone
(\S\ref{ssec:formal}(a)). (J) and (M) are published theorems, so Theorem~\ref{thm:shell} is unconditional in the usual sense; we state them as
hypotheses because they are the only inputs that we have neither re-derived nor formalised.

\begin{table}[ht]\centering\small
\caption{Lower bounds for $\liminf_{Q\to\infty}\sum_\chi N^s_{0,\chi}(T,2T)/\sum_\chi N_\chi(T,2T)$ over $\FQ$ ($1<q\le Q$) at
$T=(\log Q)^{r_0+\eps}$, $r_0\ge3$. The labels are those of \S\ref{ssec:formal}.}\label{tab:results}
\begin{tabular}{@{}>{\raggedright\arraybackslash}p{2.6cm}>{\raggedright\arraybackslash}p{3.7cm}>{\raggedright\arraybackslash}p{7.8cm}@{}}
\toprule
previous~\cite{DSouza} & this paper & status\\
\midrule
$0.7212$ & $\mathbf{0.7235}$ (Thm~\ref{thm:one-prime}) & proved on paper; in Lean \tagL\ (in the released library, \S\ref{ssec:formal})\\[3pt]
 & $\mathbf{0.9059137927}$ (Thm~\ref{thm:shell}) & proved on paper, with (J), (M) quoted \tagC; in Lean \tagL\ with (J), (M) as hypothesis (in the released library, \S\ref{ssec:formal})\\
\bottomrule
\end{tabular}
\end{table}

The constant $0.7235$ is the truncation of $2-\Bfun^{\rm K}_{C^+}(v_{\rm d})=0.7235110849\ldots$, where $\Bfun^{\rm K}$ is the functional~\eqref{eq:BK}
and $v_{\rm d}$ is the design profile of~\cite{DSouza}; $0.9059137927$ is the truncation of $2-\Bfun_{\alphap}(v_p)=0.90591379273786\ldots$ for the
degree-$12$ profile S53-L75 (\S\ref{sec:cert}). Both values are exact rationals computed in Lean's kernel. The preprint's $0.7212$ is a Lean
theorem with no hypothesis in the library \texttt{ZetaQ} of the same repository (\S\ref{ssec:formal}). For the moduli $Q/2<q\le Q$ the
preprint proved $0.7098$; the corresponding Shell value is discussed in Remark~\ref{rem:shell-dyadic} and is not claimed.

\subsection{What is new}\label{ssec:new}
The method is that of~\cite{DSouza}, which runs the rank--trace certificate of Alp\"oge and Furman~\cite{AF} on the family Gram matrix of $\FQ$
(\S\ref{sec:setting}). The certificate converts an upper bound $\frob{\widehat G}^2\le(\Bfun_*+o(1))\mathcal N$ for the Frobenius norm of the
normalised family Gram matrix into the lower bound $2-\Bfun_*$ for the proportion of simple critical zeros (Proposition~\ref{prop:3.1}).
In~\cite{DSouza}, $\Bfun_*$ is the value of the functional
\[
\Bfun(v)=\psi(0)+\int k(\alpha)\psi(\alpha)\,d\alpha,\qquad \psi=v\star v,
\]
at a window profile $v$, with kernel $k(\alpha)=\abs\alpha$ on the in-zone $\abs\alpha\le1$ and $k(\alpha)=C\abs\alpha$ on the out-zone
$1<\abs\alpha\le\lambda$, where $C=\pi^4/18$ is the loss of the multiplicative large sieve applied pointwise in the dual variable. Both theorems
of this paper change only the out-zone kernel.

\emph{The killed kernel} (\S\ref{sec:lemmaK}). For $\abs\alpha>1$ the prime sum $S(\vartheta)=\sum_na_n e(n\vartheta)$ whose large-sieve mass is
being bounded puts a fixed proportion of its mass on short arcs around the rationals $a/r$ of small denominator, and no Farey point of level $Q$
other than $a/r$ itself comes near these arcs. Gallagher's form of the large sieve with a zero set (Lemma~\ref{lem:K-G}) turns this into the
kernel $C\min(\abs\alpha,1)$ (Lemma~\ref{lem:K}), unconditionally. The design profile of~\cite{DSouza} then gives $0.7235$.

\emph{The Shell kernel} (\S\ref{sec:shell}). On an initial segment $1<\abs\alpha\le\alphap$ of the out-zone the kernel can be lowered to $1$. Assuming
the generalised Riemann hypothesis, Chandee, Lee, Liu and Radziwi\l\l~\cite[Theorem~2]{CLLR14} give an asymptotic formula, uniformly on every
compact subset of $(-2,2)$, for their $q$-analogue of the pair correlation function, averaged over the primitive characters of conductor $q\asymp Q$
with a smooth weight in $q/Q$ divided by $\varphi(q)$, and with Mellin-transform weights on the zeros: it is
$(1+o(1))\min(\abs\alpha,1)$ plus terms that vanish for each fixed $\alpha\ne0$ once $Q$ is large, so its limit is $1$ for each fixed $1<\abs\alpha<2$ [V].
Their family and their weighting of the zeros differ from ours; on $(1,\alphap]$ the kernel $1$ is the value that their theorem gives under that
hypothesis \tagM. A signed Farey identity (Lemma~\ref{lem:shell-1}) writes the family sum as a weighted sum over Farey fractions; outside
the shells around $a/r$ with $r\le R_1$ the weighted Farey density is asymptotically constant (Lemmas~\ref{lem:shell-2} and~\ref{lem:shell-P}),
inside the holes it vanishes, and the mass on the rings around the holes is expressed through the explicit formula for the primitive characters of
modulus at most $R_1$ (Proposition~\ref{prop:shell-Z}) and bounded by two zero-density estimates for the family. These estimates reach
$\alphap<5/3$ (\S\ref{ssec:shell-density}); beyond $\alphap$ the killed kernel is used. The window is taken at bandwidth $\lambda=191/100$, much
larger than the $\lambda^*=1.2507$ of~\cite{DSouza}, and the only condition on the bandwidth that the proofs use in an essential way is
$\lambda_{\rm eff}<2$ (\S\ref{ssec:bandwidth}).

\emph{Formalisation.} Theorem~\ref{thm:one-prime}, and Theorem~\ref{thm:shell} with (J) and (M) as its hypothesis, are Lean theorems in the
release of the formalisation (\S\ref{ssec:formal}); the formalisation found several statements of the
working draft false as written; the paper states the corrected forms and lists the changes in Appendix~\ref{app:corr}.

\subsection{What the results are not}\label{ssec:not}
The theorems are lower bounds for family averages. They say nothing about any individual $L(s,\chi)$: a subfamily of density zero could carry all
its zeros off the line. They say nothing about the generalised Riemann hypothesis.

Both theorems are asymptotic, and the thresholds in the proofs are very large. For Theorem~\ref{thm:shell}, the condition $X\le Q^2$ that makes the
large sieve lossless starts near $\log_{10}Q\approx177$ at $r_0+\eps=3.5$ (\S\ref{ssec:bandwidth}); the existence of the design and the moment floor
$a\ge3/4$ need $\log Q$ at least of order $10^5$, and the formal proof's threshold is larger and implicit (\S\ref{sec:design}); the error term
$(\log Q)^{-\theta}$, $\theta<503/1994$, decreases slowly. For Theorem~\ref{thm:one-prime}, the prime number theorem enters Lemma~\ref{lem:K-P}
through the error term of $\sum_{n\le y}\Lambda(n)$ at $y\asymp QT\Lc$ with a fixed power of $\log Q$ to spare; with the explicit constants
of~\cite{FKS} the corresponding row first becomes small near $\log Q\approx5900$ (Remark~\ref{rem:K-asymptotic}). No finite-$Q$ statement
of~\cite{DSouza} is improved by either theorem, and none is claimed.

Theorem~\ref{thm:shell} rests on the two zero-density estimates (J) and (M), which we quote from the literature; their proofs are long and are not
formalised (\S\ref{ssec:shell-density}). Its rate is much slower than the rate $\log\log Q/\log Q$ of Theorem~\ref{thm:one-prime}.

We do not claim the Shell constant for the moduli $Q/2<q\le Q$ (Remark~\ref{rem:shell-dyadic}), a bounded-height version of either theorem, the
corresponding statements for the even and odd subfamilies, or any statement about distinct zeros; the working draft contains arguments for these
(\S\ref{sec:remarks}), and they have not been brought to the standard of this paper.

\subsection{Relation to other work}\label{ssec:others}
Statements about other work carry the tag [V] when we read them in the source, \Vsim\ when we read them only through a description (the authors'
own, or that of the preprint~\cite{DSouza}), and \tagM\ when they are our inference.

\emph{The source of the method.} Alp\"oge and Furman~\cite{AF} prove unconditionally that at least two thirds of the nontrivial zeros of $\zeta$,
counted with multiplicity, are simple and lie on the critical line, and at least five sixths are distinct; with the Montgomery--Taylor window the
constants become $0.6725\ldots$ and $0.8362\ldots$ (their Theorem~A). Their Theorem~B states that Theorem~A holds verbatim for $L(s,\chi)$ for any
fixed primitive Dirichlet character $\chi$, and their abstract states that the results are formally verified in Lean~4 [V]. The rank--trace
certificate, the Gram realisation of the explicit formula and the other inputs listed as H1--H6 in \S\ref{ssec:H} are theirs. In their
Remark~7.2 they record, with the details omitted, that summing over primitive $\chi\bmod q$, $q\le Q$, with the smooth weights $(1-q/Q)^2$, a
weighted large-sieve inequality and a Gevrey taper give family-average proportions of at least $0.811$ simple and on the line and $0.905$
distinct, unconditionally, for $T$ a power of $\log Q$; they state that this has been checked by independent derivations and is not included in
their Lean formalisation [V]. That family has a different weight from ours, so we do not compare the numbers \tagM. The family version for the
unweighted family $1<q\le Q$ at polylogarithmic height, with its uniformity in $q$, is the work of~\cite{DSouza}.

\emph{Other family averages.} Conrey, Iwaniec and Soundararajan~\cite{CIS2} use their asymptotic large sieve~\cite{CIS1} and Levinson's method to
prove that at least $14/25$ of the zeros with $\abs\gamma\le T$, $(\log Q)^6\le T\le(\log Q)^A$, are simple and on the critical line, on average over
the primitive characters with the weight $\Psi(q/Q)/\varphi(q)$, $\Psi$ smooth with compact support in $(0,\infty)$, so that $q\asymp Q$ [V].
Sono~\cite{Sono21}, ``instead of the asymptotic large sieve'', uses a formula of Conrey, Iwaniec and Soundararajan for the mean square of the product
of a Dirichlet $L$-function and a Dirichlet polynomial, with shifts~\cite{CIS19} (by its abstract another application of their work on the asymptotic
large sieve), which he extends to shifts with imaginary parts up to a power of $\log Q$, together with Feng's mollifier; he proves that at least $0.6044$ of the zeros with $T\le\gamma<2T$, $(\log Q)^2\le T\le(\log Q)^A$, are
simple and on the critical line, on average over the even primitive characters with a smooth weight in $q/Q$ supported in $[1,2]$ [V]. Under the
generalised Riemann hypothesis, Chandee, Lee, Liu and Radziwi\l\l~\cite{CLLR14} study low-lying zeros with the asymptotic large sieve, and
Sono~\cite{Sono16} uses their pair correlation formula [V]. This paper uses none of these inputs. Our family is the unweighted family of all primitive
characters with $1<q\le Q$; the families of~\cite{CIS2}, \cite{Sono21} and~\cite{CLLR14} carry smooth conductor weights with $q\asymp Q$, that
of~\cite{Sono21} contains the even characters only, \cite{CIS2} counts the window $\abs\gamma\le T$, and~\cite{CLLR14} and~\cite{Sono16} count
low-lying zeros under that hypothesis; so we do not compare the numbers \tagM. (We read the arXiv versions
of~\cite{CIS2}, \cite{Sono21} and~\cite{CLLR14}, and the publisher's abstract of~\cite{Sono16}.) The preprint~\cite{DSouza} reports that it found no
earlier result for the unweighted family of all primitive characters of modulus $q\le Q$ \Vsim; we have not searched further.

\emph{Per-character claims.} In unrefereed Zenodo records, Hongyi Yang and Shihua Yang state bounds for $\zeta$ which, according to the records,
hold for every fixed primitive Dirichlet $L$-function as well: at least $0.7962$ of the zeros simple and on the critical line and $0.8981$ distinct,
with ineffective constants, a result the authors declare ``at the certified-candidate level, not as an established theorem''~\cite{YY1}; and, in a
second record~\cite{YY2}, that almost all zeros are simple and on the critical line, where the record states that the constant layers are proved,
that the analytic chain is a certified candidate pending external review, and that the certificates are kernel-checked in Lean~4; the record page
shows its files under embargo, with the stated reason that the version contains errors and needs revision [V] (record pages as accessed on
28~September~2026). These are statements for each fixed character as the height tends to infinity; they are not averages over $q\le Q$ at
polylogarithmic height, and neither kind of statement implies the other \tagM. We have not verified these claims.

\subsection{Formal verification and status}\label{ssec:formal}
A Lean~4 formalisation~\cite{Lean4} accompanies this paper (toolchain v4.33.0-rc2, Mathlib~\cite{Mathlib} revision 51e6992). It is built on the
Lean artifact \texttt{Zeta23} of~\cite{AF} and on the library \texttt{ZetaQ} that formalises~\cite{DSouza}, and it is released in the repository
\repourl, at tag \repotag, whose Lean sources are those of commit \relcommit\ of the development history (the repository is a snapshot of the
release). The family library of this paper is its library \texttt{ZetaShell}, and Appendix~\ref{app:lean-release} gives the commands that check
it. A declaration is \emph{kernel-checked} if \texttt{\#print axioms} reports only \texttt{propext}, \texttt{Classical.choice} and
\texttt{Quot.sound}: no \texttt{sorry}, no added axiom, no \texttt{native\_decide}. In the release the library \texttt{ZetaShell} has $327$ files,
$51{,}318$ lines and $1{,}744$ public theorems; $1{,}708$ of them are kernel-checked, and $36$ depend on \texttt{sorryAx}: $29$ open statements,
each with its own \texttt{sorry}, and $7$ theorems derived from them. No other axiom and no \texttt{native\_decide} occurs. The formal statuses
below are those of the census of the release, one Lean job, run on \censusref, that prints, for each public theorem, its axioms and the open
statements its proof reaches (its output is in \texttt{audit/followup/} of the repository); Appendix~\ref{app:lean} quotes it
(Appendix~\ref{app:lean-census}).

We mark a theorem quoted from the literature by \tagC\ and a numerical computation outside Lean by \tagN. For the formal status we use the
following labels.

\begin{enumerate}[label=(\alph*)]
\item \tagL\ \emph{Kernel-checked.} The only hypothesis that is not part of the mathematical statement is \texttt{ZeroDensityInput}, that is (J) and
(M) (see below), and it occurs exactly where the paper assumes (J) and (M).
In the library \texttt{ZetaQ}: the theorems of~\cite{DSouza} at $0.7212$ and $0.7098$, for every height exponent $\ge3$
(\texttt{ZetaQ.JoinProved.theorem\_one\_generic\_proved'} and \texttt{corollary\_two\_dyadic\_proved'}), since the large-sieve input was re-proved in
Gallagher's form.
Theorem~\ref{thm:one-prime}, \texttt{ZetaShell.LemmaK.theorem\_one\_prime\_design}, in the library \texttt{ZetaShell} of the release: in the census it
depends on the axioms \texttt{propext}, \texttt{Classical.choice} and \texttt{Quot.sound} only and rests on no open statement, and the file
\texttt{comparator/PrintAxioms/FollowUpShell.lean} prints the same axioms (Appendix~\ref{app:lean-release}). Its hypotheses are $3\le r$ and $0<\eps$
(here $r$ is our $r_0$), and its constant is $7235/10000$. The companion statement at $0.7237$ (\texttt{theorem\_one\_prime}) still rests on one open
node.
Also in the library \texttt{ZetaShell}: the certificates \texttt{ZetaShell.certS53} and \texttt{certD53}; the form of the prime number theorem that
the proofs use, $\sum_{n\le x}\Lambda(n)-x=o(x(\log x)^{-A})$ for every $A$; the signed Farey identity, Lemma~\ref{lem:shell-6b},
Lemma~\ref{lem:shell-F1} and Lemma~\ref{lem:shell-1prime}(a)--(c); Lemma~\ref{lem:shell-2}(b) and~(c); the design layer at $\lambda=191/100$ except its
zone-slope node; Lemmas~\ref{lem:K-farey}, \ref{lem:K-G} (in corrected form), \ref{lem:K-Z}, \ref{lem:K-R}, \ref{lem:K-S} and
$\sum_{r\le R}\mu^2(r)/\varphi(r)\ge\log(R+1)$; the hole bounds~\eqref{eq:K-hole}; Lemma~\ref{lem:shell-5a}, the Weil form of
Lemma~\ref{lem:shell-5b} and Lemma~\ref{lem:shell-5c}(2),(2$'$); Lemma~\ref{lem:shell-6c}; Proposition~\ref{prop:shell-Z} in the Weil form
(\texttt{propZ\_W}); Lemma~\ref{lem:shell-3} (\texttt{signed\_gallagher}); Lemma~\ref{lem:K}(i) and Lemma~\ref{lem:K-P}; the Frobenius-row chain at the
new bandwidth (\texttt{F1c\_chain}); the strip, the integrability, the pointwise part of the killed zone and the majorant construction of
Theorem~\ref{thm:shell-K}; the family errors and the elementary terms of Theorem~\ref{thm:shell-S}. The individual statements and their status are in
Appendix~\ref{app:lean}.

The statement of record of Theorem~\ref{thm:shell}, \texttt{ZetaShell.shell\_S53\_qle\_reduced}, depends on the three standard axioms only and rests
on no open statement. It is stated against counting functions defined from Mathlib alone and proved equal to those of \texttt{ZetaQ}, with the single
displayed hypothesis \texttt{ZeroDensityInput}, which is (J) and (M) as stated in the sources (Appendix~\ref{app:lean-record}). In the release both
theorems of this paper are also stated in the trusted files \texttt{comparator/ChallengeDeps/FollowUpFamily.lean}, which imports Mathlib alone, and
\texttt{comparator/Challenge/FollowUpFamily.lean}, as \texttt{family\_simple\_on\_line\_killed} (Theorem~\ref{thm:one-prime}, no hypothesis) and
\texttt{family\_simple\_on\_line\_shell} (Theorem~\ref{thm:shell}, with \texttt{ZeroDensityInput} as its only hypothesis); these files are the place to
read what is claimed. Both are proved in \texttt{comparator/Solution/FollowUpFamily.lean} from \texttt{theorem\_one\_prime\_design} and
\texttt{shell\_S53\_qle\_reduced}, and \texttt{comparator/PrintAxioms/FollowUpFamily.lean} prints both at the three standard axioms
(Appendix~\ref{app:lean-release}); the comparator tool itself has not been run on them. The statement of record is proved through the frame with the
Frobenius row derived from Theorem~\ref{thm:shell-K}, in a form that does not use the strong prime number theorem \texttt{PNTErrorTerm}
(\texttt{shell\_frame\_qle\_of\_K\_noP}; the form \texttt{shell\_frame\_qle\_of\_K}, which also assumes \texttt{PNTErrorTerm} and \texttt{CertS53}, is
kernel-checked as well). Below it, also kernel-checked: Lemma~\ref{lem:shell-star} in its corrected form
(\texttt{ring\_le\_primitive\_corr}, with the Gauss sums of imprimitive characters, which Mathlib lacks, in a new module); Theorem~\ref{thm:shell-S}
(\texttt{shell\_S}) with its steps S1 and S2, Lemma~\ref{lem:shell-6a} (\texttt{Z6a\_zero\_sums}) and Lemma~\ref{lem:shell-6d} in the corrected form
\texttt{Z6d\_counts\_corr}; Lemma~\ref{lem:shell-2}(a) (\texttt{lemma2a}) and Lemma~\ref{lem:shell-P} (\texttt{lemmaP}), with \texttt{line\_sum\_asymp},
\texttt{s3\_tail\_plain\_pt} and \texttt{s3\_tail\_sum}; and Theorem~\ref{thm:shell-K} (\texttt{thmK}) with its integrated killed-zone part
(\texttt{K3b\_integrated}), its Shell-zone part (\texttt{K2\_shellZone}) and the pointwise
assembly~\eqref{eq:shell-assembly} in the corrected form \texttt{AS\_pointwise\_corr}; the transfer (\texttt{KT\_transfer}) is kernel-checked there too,
but the proof of the Shell-zone part does not use it. The corrected forms \texttt{Z6d\_counts\_corr} and \texttt{AS\_pointwise\_corr} each add one
hypothesis to the first formal form of an internal statement (Appendix~\ref{app:corr}); the first forms are kept in the library, open and unused.
\item \tagLm\ \emph{Kernel-checked modulo named open nodes.} The frame of the formal Shell statement with the Frobenius row as one node
(\texttt{shell\_frame\_qle\_of\_nodes}) rests on one open node, \texttt{frob\_row\_shell}, which is the Frobenius row of
Theorem~\ref{thm:shell}, that is, the Shell argument. The proof of the statement of record does not use it: it derives the Frobenius row from
Theorem~\ref{thm:shell-K} and the Frobenius-row chain (see (a)).
\item \tagLs\ \emph{Stated in Lean, not proved.} Of the statements this paper names: the first formal forms of Lemma~\ref{lem:shell-6d}
(\texttt{Z6d\_counts}) and of the pointwise assembly (\texttt{AS\_pointwise}), which the corrected forms replace and which no proof uses
(Appendix~\ref{app:corr}); the Frobenius row \texttt{frob\_row\_shell} of (b); the zone-slope statement of \S\ref{sec:design}(viii), which the proof
of the statement of record does not use; and earlier forms of the headline and the dyadic and bounded-height versions of the statement of record,
which are not claimed (Appendix~\ref{app:lean-record}). None of these, and no other open statement of the census, is reached by the proof of either
theorem (Appendix~\ref{app:lean-census}).
\item \emph{Not formalised at all.} The zero-density estimates themselves, which enter only as the displayed hypothesis; the asymptotic consequences of
Lemma~\ref{lem:shell-1prime}; Lemma~\ref{lem:shell-4} in the form used by the Shell route (for its analogue inside Lemma~\ref{lem:K} see (a)). In the
release the proof of the pointwise assembly proves, in its own form, the bounds that it needs from Lemmas~\ref{lem:shell-1prime} and~\ref{lem:shell-4}.
\end{enumerate}
The formalisation found statements of the working draft false as written; the ones that concern this paper are listed in Appendix~\ref{app:corr}. In
each case the corrected statement holds wherever the statement is used, and no constant of Theorems~\ref{thm:one-prime} or~\ref{thm:shell} changed; the
constant $0.72375$ that the working draft stated for Theorem~\ref{thm:one-prime} is not reached inside the formal design, and this paper states $0.7235$
(\S\ref{ssec:K-constants}).

\subsection{How the work was done}\label{ssec:how}
\emph{Disclosure of generative-AI use.} This paper and its Lean development were produced substantially by a generative AI system, Claude
(Anthropic), working under the direction of the named author, who takes full responsibility for its contents. No AI system is listed as an author.

The arguments, the proofs on paper and the Lean proofs were written by instances of Claude working as separate agents on written tasks, which also
checked one another's work; the formal statuses are those that Lean reports (\S\ref{ssec:formal}). The \emph{working draft} to which this paper refers
is the project's unpublished draft of about $320$ pages, which contains, besides the material of this paper, further families and conditional
programmes; where this paper departs from it or from~\cite{DSouza}, the change is listed in Appendix~\ref{app:corr}.

\emph{Scope of the verification.} Theorems~\ref{thm:one-prime} and~\ref{thm:shell} are Lean theorems that depend only on the three standard
axioms. Theorem~\ref{thm:one-prime} has no hypothesis; Theorem~\ref{thm:shell} assumes the zero-density estimates \textup{(J)} and
\textup{(M)}, which are published theorems cited here and not proved. Some intermediate error rates of the proofs on paper are not checked in
Lean, because the formal proof absorbs those errors into the slack of the certificate (\S\ref{sec:proof}). The paper has not been refereed,
and the Lean development has not been rebuilt outside this project.

\subsection{Plan}
Section~\ref{sec:setting} recalls the setting of~\cite{DSouza}, the notation, the imported hypotheses and the role of the bandwidth.
Section~\ref{sec:lemmaK} proves the killed-kernel lemma and Theorem~\ref{thm:one-prime}. Section~\ref{sec:shell} is the Shell route,
Section~\ref{sec:design} the design at the new bandwidth, Section~\ref{sec:cert} the certificate and Section~\ref{sec:proof} the proof of
Theorem~\ref{thm:shell}. Section~\ref{sec:remarks} discusses rates, thresholds and what would improve the constants. Appendix~\ref{app:lean} reproduces
the Lean statements; Appendix~\ref{app:corr} lists the corrections to the working draft, to formal statements and to~\cite{DSouza}.
